\section*{Заключение}

В процессе выполнения НИР были изучены принципы организации
виртуальной вычислительной среды, имитирующей ресурсы вычислительного
кластера, а также технологии управления заданиями средствами системы
SLURM на базе дистрибутива Linux Fedora Server в среде виртуализации
Oracle VM VirtualBox. В результате работы были выполнены поставленные
задачи:

\begin{enumerate}
\item Установлена и настроена среда виртуализации Oracle VM VirtualBox.

\item Создано пять виртуальных машин: один управляющий узел и четыре
вычислительных узла кластера.

\item Установлен дистрибутив Linux Fedora Server (архитектура aarch64)
на все виртуальные машины в режиме интерфейса командной строки.

\item Выявлена и устранена проблема совместимости эмулируемого
контроллера накопителей с архитектурой ARM хост-системы путём перехода
на контроллер VirtIO SCSI.

\item Настроена виртуальная локальная сеть между узлами кластера с
разрешением имён посредством файла \texttt{/etc/hosts}.

\item Настроено подключение по протоколу SSH с авторизацией по
криптографическим ключам между управляющим и вычислительными узлами.

\item Установлена и настроена служба аутентификации munge, а также
диспетчер заданий SLURM (\texttt{slurmctld} на управляющем узле и
\texttt{slurmd} на вычислительных узлах).

\item Выполнено клонирование настроенного вычислительного узла для
создания оставшихся узлов кластера с последующей индивидуализацией их
сетевых параметров.

\item Проведён анализ набора учётных данных \texttt{acct\_0921-0923}
по заданным атрибутам варианта (UID~50430, JobName~``run\_bwa.sh''),
вычислены границы диапазона фактического времени выполнения задач, а
также параметры $\mu$ и $\sigma$ логарифма ошибки оценки времени
выполнения.

\item Разработан имитатор пользовательских заявок, генерирующий
фактическое время выполнения задачи по нормальному распределению в
вычисленном диапазоне и пользовательскую оценку времени выполнения по
логнормальному распределению с вычисленными параметрами, с
автоматической постановкой сформированных заявок в очередь диспетчера
SLURM.

\item Реализован сбор и сохранение в файл статистики поставленных в
очередь заявок для последующего сравнительного анализа плановых и
фактических характеристик их выполнения.
\end{enumerate}

Таким образом, в ходе выполнения НИР развёрнута работоспособная
виртуальная модель вычислительного кластера, на базе которой
реализован и экспериментально проверен имитатор пользовательской
нагрузки, статистические характеристики которого согласуются с
параметрами, полученными в результате анализа реального набора учётных
данных. Полученные результаты могут быть использованы в качестве
основы для дальнейшего расширения функциональности диспетчера и более
детального анализа эффективности алгоритмов планирования заданий.